\emph{Set-up.} Units $\hbar = c = 1$. On one generator, $\mathfrak{h} = L^{2}(\R, du)\otimes K$, $\CAR(\mathfrak{h})$ the continuum mode space (Definition \claimref{obj.hawking.continuum-mode-space}) with $a(f)$ antilinear in $f$, $\mathfrak{A}(O) = \CAR(L^{2}(O)\otimes K)$ for an interval or half-line $O$, $\tau_{t}(a(f)) = a(f(\cdot - t))$, $P = i\,d/du$ (so $e^{-ipu}$ has $P$-eigenvalue $p$ and $\tau_{t}$ acts on it by $e^{itp}$), and a gauge-invariant quasi-free state $\omega_{E}$ is fixed by its symbol $E$, $\omega_{E}(a^{*}(f)a(g)) = \ip{g}{Ef}$. The null vacuum is $\omega_{0} = \omega_{E_{0}}$, $E_{0} = \chi_{(-\infty,0)}(P)\otimes\id_{K}$, kernel $\tfrac{i}{2\pi}(u - v + i0)^{-1}$; it is the unique gauge-invariant quasi-free state of (N$_{+}$), whose ground-state clause asks that $\tau$ be implemented in the GNS representation by a unitary group with nonnegative generator fixing the GNS vector (Theorem \claimref{synthesis.positive-null-energy-gives-r-premises}(i)). Its half-line modular data are geometric: the modular group of $\mathfrak{A}((a, \infty))$ is the boost about $a$ at rapidity $2\pi t$, and the modular conjugation $J^{0}_{a}$ acts as the reflection $r_{a}(u) = 2a - u$ composed with charge conjugation: $J^{0}_{a}\pi_{0}(a(f))J^{0}_{a} = \Theta\,\pi_{0}(a^{*}(c\,f\circ r_{a}))$ for $f$ supported in $(a, \infty)$, with $\Theta$ the grading unitary normalised by $\Theta\Omega_{0} = \Omega_{0}$ and $c$ a constant phase (Fact \claimref{angle.imported-theorems-deliver-their-hypotheses}: Bisognano-Wichmann and the conformal-net form, the modular conjugation being CPT composed with the turn by $\pi$; its action on the fields as the twisted reflection, Fact \claimref{synthesis.imported-modular-conjugation-of-the-chiral-fermi-field-is-the-twisted-reflection}(c)). In the Klein-twist form $J^{0}_{a}\psi(f)J^{0}_{a} = \eta Z\psi^{*}(\rho_{a}f)Z^{*}$, $Z = (1 + i\Theta)/(1 + i)$, one has $Z\psi^{*}Z^{*} = i\Theta\psi^{*}$, so $c = i\eta$; the companion field-level result \claimref{synthesis.geometric-modular-conjugations-on-a-generator-give-the-null-ground-state} (Step 4, natural-cone positivity, proved and checked there) fixes $\eta = -1$ for the null vacuum, so $c = -i$. Steps 1, 2, 4 and 5 use only that $c$ does not depend on $f$. For an automorphism $\alpha$ the state $\omega\circ\alpha$; for an antilinear automorphism $\theta$ the state $x \mapsto \overline{\omega(\theta(x))}$.

\emph{Step 1: net-level identity, (i).} Let $m : \R \to U(K)$ be measurable and $\alpha_{m}(a(f)) = a(mf)$; multiplication by $m$ is unitary and preserves the anticommutator, so $\alpha_{m}$ is an automorphism, and since multiplication by $m$ and by $m^{-1}$ maps $L^{2}(O)\otimes K$ onto itself, $\alpha_{m}(\mathfrak{A}(O)) = \mathfrak{A}(O)$ for every $O$. The triple $(\mathcal{H}_{0}, \pi_{0}\circ\alpha_{m}, \Omega_{0})$ represents $\omega_{0}\circ\alpha_{m}$, and $\Omega_{0}$ is cyclic for it because $\alpha_{m}$ is onto; by uniqueness of the GNS representation it is the GNS triple. Hence $\pi_{0}(\alpha_{m}(\mathfrak{A}(O)))'' = \pi_{0}(\mathfrak{A}(O))''$ for every $O$: the two states define the same Hilbert space, vector and von Neumann algebra for every region. The modular operator and conjugation of a pair (von Neumann algebra, cyclic separating vector) depend on nothing else, so they coincide region by region, and so does every property of the net and the vector built from them: the generator form of the Condition of Geometric Modular Action (in the Buchholz-Summers form, a condition on the algebras as sets, Fact \claimref{synthesis.imported-buchholz-florig-summers-geometric-modular-action-and-modular-stability}(a) \cite{BuchholzFlorigSummers1999, BuchholzSummers1993}), geometric action of the modular groups on the net, modular stability, twisted Haag duality, and the vacuum's translation unitaries $U_{0}(t)$, with nonnegative generator, which map $\pi_{0}(\mathfrak{A}(O))''$ onto $\pi_{0}(\mathfrak{A}(O + t))''$ and fix $\Omega_{0}$.

\emph{Step 2: the Fermi seas, (ii).} Take $m(u) = e^{iu\kappa}$, $\kappa = \kappa^{*}$ on $K$. Then $\omega_{0}(\alpha_{m}(a^{*}(f)a(g))) = \ip{mg}{E_{0}mf} = \ip{g}{m^{-1}E_{0}mf}$, and $m^{-1}$ followed by $P$ and then $m$ is $P - \kappa$ on plane waves ($m$ sends the eigenvalue $p$ to $p - \kappa$), so the symbol is $E_{\kappa} = \chi_{(-\infty,0)}(P - \kappa)$, a projection: $\omega_{\kappa}$ is pure, and it is $\tau$-invariant since $E_{\kappa}$ commutes with $P$. On a bounded interval $O$ of length $L$, for $\kappa = k\,\id_{K}$ with $k > 0$ (the general case is the sum over eigenvectors of $\kappa$, with the sign of each eigenvalue), $\chi_{O}(E_{k} - E_{0})\chi_{O} = \frac{1}{2\pi}\int_{0}^{k}\chi_{O}\ket{e_{p}}\bra{e_{p}}\chi_{O}\,dp$, $e_{p}(u) = e^{-ipu}$, a positive operator of trace $kL\dim K/2\pi$: trace class, so by the Powers-Stormer inequality and criterion (Fact \claimref{synthesis.imported-powers-stormer-takesaki-disjointness}(a), (c)) the restrictions to $\mathfrak{A}(O)$ are quasi-equivalent, the same argument as in the stored local-normality theorem. The kernel difference $\frac{1}{2\pi}\int_{0}^{k}e^{-ip(u - v)}\,dp$ is entire in $u - v$, so the generator microlocal condition holds at every point (\claimref{synthesis.microlocal-condition-on-a-generator-gives-the-local-null-ground-state}, which characterises it by a smooth kernel difference).
Not a ground state: since $\omega_{\kappa}$ is $\tau$-invariant, $\tau$ has a canonical implementer $U_{\kappa}(t)$ fixing $\Omega$, and $U_{\kappa}(t)\pi_{\kappa}(a^{*}(f))\Omega = \pi_{\kappa}(a^{*}(e^{itP}f))\Omega$, $U_{\kappa}(t)\pi_{\kappa}(a(g))\Omega = \pi_{\kappa}(a(e^{itP}g))\Omega$. For $f$ in the range of $1 - E_{\kappa}$ (a particle at momentum $p > k$) the energy is $+p$; for $g$ in the range of $E_{\kappa}$ (a hole at $p < k$), $a$ being antilinear, it is $-p$. With $k > 0$, holes at $p \in (0, k)$ have negative energy, and $n$ holes in disjoint momentum windows inside $(k/2, k)$ have energy below $-nk/2$; for $k < 0$ the particles at $p \in (k, 0)$ do the same. In an irreducible representation any other implementer differs by a character $e^{i\lambda t}$, so no implementer has a generator bounded below.
On the gauge-invariant subalgebra $\mathfrak{A}^{G}$: $\alpha_{m}\tau_{t}\alpha_{m}^{-1}(a(f)) = a(e^{iu\kappa}e^{-i(u - t)\kappa}f(u - t)) = a(e^{it\kappa}f(\cdot - t))$, $\tau_{t}$ followed by the internal rotation $e^{it\kappa}$; inversely $\alpha_{m}^{-1}\tau_{t}\alpha_{m}(a(f)) = a(e^{-it\kappa}f(\cdot - t))$, whose one-particle generator on $e^{-ipu}$ is $p - \kappa$, and $\omega_{\kappa} = \omega_{0}\circ\alpha_{m}$ is the ground state of this lift ($\omega_{\kappa}\circ\alpha_{m}^{-1}\tau_{t}\alpha_{m} = \omega_{0}\circ\tau_{t}\circ\alpha_{m}$, implemented by $U_{0}$). The rotation $e^{it\kappa}$ commutes with every gauge rotation; for $\kappa = k\,\id$ it is a gauge rotation and acts trivially on $\mathfrak{A}^{G}$, so $\alpha_{m}$ and $\tau_{t}$ commute there. Then $U_{0}(t)\pi_{0}(\alpha_{m}(x))U_{0}(t)^{*} = \pi_{0}(\tau_{t}\alpha_{m}(x)) = \pi_{0}(\alpha_{m}(\tau_{t}x))$ for $x \in \mathfrak{A}^{G}$: on the cyclic subspace $\overline{\pi_{0}(\mathfrak{A}^{G})\Omega_{0}}$, which $U_{0}$ preserves, $U_{0}$ implements $\tau$ for $\omega_{\kappa}$ with nonnegative generator and fixes the vector, so $\omega_{\kappa}$ restricted to $\mathfrak{A}^{G}$ is a ground state for $\tau$. It differs from $\omega_{0}$ there: $\omega_{k}(a^{*}(f)a(f)) - \omega_{0}(a^{*}(f)a(f)) = \ip{f}{(E_{k} - E_{0})f}$, a charge density $k/2\pi$, and the energy density is $\frac{1}{2\pi}\int_{0}^{k}p\,dp = k^{2}/4\pi$ per dimension of $K$. With $K$ two-dimensional and $\kappa$ of eigenvalues $\pm k$ the charges cancel and the energies add; this state is not a ground state for $\tau$ (the eigenvalue $k$ gives negative-energy holes as above), and it is the ground state of the lift $\alpha_{m}^{-1}\tau_{t}\alpha_{m}$, $\tau_{t}$ followed by $e^{-it\kappa}$ (one-particle generator $P - \kappa$).
Elementwise conjugation: with $\pi_{\kappa} = \pi_{0}\circ\alpha_{m}$ and $J_{a} = J^{0}_{a}$ (Step 1), $J_{a}\pi_{\kappa}(a(f))J_{a} = J^{0}_{a}\pi_{0}(a(mf))J^{0}_{a} = \Theta\,\pi_{0}(a^{*}(c\,(mf)\circ r_{a})) = \Theta\,\pi_{\kappa}(a^{*}(c\,m^{-1}\,(mf)\circ r_{a}))$, and $m(u)^{-1}m(2a - u) = e^{2i\kappa(a - u)}$, constant in $u$ only for $\kappa = 0$.

\emph{Step 3: the sign, (iii).} $R : u \mapsto -u$ is unitary on $\mathfrak{h}$ with $RPR = -P$, so $\omega_{0}\circ\alpha_{R}$ has symbol $RE_{0}R = \chi_{(0,\infty)}(P)$ (the point $0$ is null), the filled sea $\bar\omega_{0}$. With $\phi(f) = a(f) + a^{*}(f)$, real-linear in $f$, $a(f) = \tfrac12(\phi(f) + i\phi(if))$; the antilinear automorphism $\theta$ fixing every $\phi(f)$ has $\theta(a(f)) = \tfrac12(\phi(f) - i\phi(if)) = a^{*}(f)$, and $x \mapsto \overline{\omega_{0}(\theta(x))}$ has $\overline{\omega_{0}(a(f)a^{*}(g))} = \overline{\ip{f}{(1 - E_{0})g}} = \ip{g}{(1 - E_{0})f}$: symbol $1 - E_{0} = \chi_{(0,\infty)}(P)$ again. $\theta$ commutes with $\tau$ ($\tau_{t}\phi(f) = \phi(f(\cdot - t))$), with the grading, and with the gauge rotations ($\gamma_{\vartheta}\phi(f) = \phi(e^{i\vartheta}f)$), and maps the real algebra generated by $\{\phi(f) : f \in L^{2}(O)\otimes K\}$ onto itself. The GNS triple of $\bar\omega_{0}$ is $(\mathcal{H}_{0}, \pi_{0}\circ\alpha_{R}, \Omega_{0})$, and the algebra of $(a, \infty)$ is $\pi_{0}(\mathfrak{A}((-\infty, -a)))''$, the twisted commutant $Z\mathcal{R}'Z^{*}$ of the vacuum's algebra $\mathcal{R}$ of $(-a, \infty)$; since $Z\Omega_{0} = \Omega_{0}$, $J^{0}_{-a}ZJ^{0}_{-a} = Z^{*}$ and $Z^{2} = \Theta$, its modular conjugation is $ZJ^{0}_{-a}Z^{*} = Z^{2}J^{0}_{-a} = \Theta J^{0}_{-a}$. Inverting the vacuum relation $J a(f) J = \Theta a^{*}(c\,f\circ r)$ ($J = J^{0}_{-a}$, $r = r_{-a}$, $J\Theta J = \Theta$): applying $J$ gives $Ja^{*}(c\,f\circ r)J = \Theta a(f)$; for $g$ supported in $(-\infty, -a)$ and $f = \bar c\,g\circ r$, $Ja^{*}(g)J = \Theta a(\bar c\,g\circ r) = c\,\Theta a(g\circ r)$, and the adjoint is $J a(g) J = \bar c\,a^{*}(g\circ r)\Theta = -\bar c\,\Theta a^{*}(g\circ r)$. Hence $(\Theta J)a(g)(\Theta J)^{-1} = \Theta(-\bar c\,\Theta a^{*}(g\circ r))\Theta = \bar c\,\Theta a^{*}(g\circ r)$. Transporting by $R$ ($R\,r_{-a}R = r_{a}$): the filled sea's conjugations act as the reflections about the cuts with the phase $\bar c$, on the same half-line algebras with the same inclusions; with $c = -i$ (set-up) this is $+i$. Consistency of the phases (argued): by (vi) in its continuum form ($\Gamma_{\Omega} \to \Theta$; the continuum one-particle polar decomposition exists since the cross-cut block is injective with dense range), $J_{0}$ acts through the polar part of $X = \chi_{(-\infty,0)}E\chi_{(0,\infty)}$; $(XF)(-x) = \int_{0}^{\infty}h(x + y)F(y)\,dy$ with $h(t) = \frac{1}{2\pi}\int_{S}e^{ipt}\,dp$ for $E = \chi_{S}(P)$, which is $-\frac{i}{2\pi t}$ for $S = (-\infty, 0)$ and $+\frac{i}{2\pi t}$ for $S = (0, \infty)$; $\bar c\,h \geq 0$ (the Carleman kernel $1/(x + y)$ is a positive operator) gives $c = -i$ for $\omega_{0}$ and $c = +i$ for $\bar\omega_{0}$, in agreement.

\emph{Step 4: independence, (iv).} $\alpha_{m} \in G_{1}$ by Step 1 and it commutes with the gauge rotations and the grading; $\theta \in G_{2}$ by Step 3. If a property is $G_{1}$-invariant and $\omega_{0}$ has it, $\omega_{0}\circ\alpha_{m}$ has it, in particular every $\omega_{\kappa}$, which for $\kappa \neq 0$ is not a ground state for $\tau$ (Step 2). If it is $G_{2}$-invariant, $\bar\omega_{0}$ has it, and in the GNS space of $\bar\omega_{0}$ the implementer of $\tau$ is $U_{0}(-t)$ transported by $R$, with nonpositive, nonzero generator, so $\bar\omega_{0}$ is not a ground state for $\tau$. The ground-state clause is not $G_{1}$-invariant because $G_{1}$ does not normalise $\{\tau_{t}\}$: $\alpha_{m}\tau_{t}\alpha_{m}^{-1}$ maps $\mathfrak{A}(O)$ onto $\mathfrak{A}(O + t)$ for every $m$, and for $m = e^{iu\kappa}$ it is $\tau_{t}$ followed by $e^{it\kappa}$ (Step 2), so fixing which lift is the translation is fixing $\kappa$, the flat connection $\kappa\,du$ along the generator. It is not $G_{2}$-invariant because positivity of a generator $A$ with $U(t) = e^{itA}$ refers to $i$: under $\theta$ the implementer's generator changes sign. Geometric algebra (identities exact; their reading argued, not proved): in $\Cl(1,3)$ with $n = \gamma_{0} + \gamma_{3}$, $B = \gamma_{0}\gamma_{3}$, $J = \gamma_{1}\gamma_{2}$, $I_{4} = \gamma_{0}\gamma_{1}\gamma_{2}\gamma_{3}$: $I_{4}B = J$, $J^{2} = -1$, $B^{2} = 1$, $JB = BJ$. The rotor $R_{\varphi} = \cos\tfrac{\varphi}{2} + \sin\tfrac{\varphi}{2}J$ (checked at $\cos\tfrac{\varphi}{2} = \tfrac35$) has $R_{\varphi}nR_{\varphi}^{-1} = n$, $R_{\varphi}BR_{\varphi}^{-1} = B$ and turns $\gamma_{1}$: a transverse rotation at a rate along the generator is invisible to $n$ and the boost, and if the field's $i$ is read as $J$, a transverse rotation at a rate along the generator acts on the spinor values as a $u$-dependent phase of the kind $\alpha_{m}$ (a reading). The transverse reflection $v \mapsto -\gamma_{1}v\gamma_{1}^{-1}$ fixes $n$ and $B$ and sends $J \mapsto -J$, $I_{4} \mapsto -I_{4}$: under the same reading, $i = J = I_{4}B$, the sign of $i$ relative to the time orientation (fixed by $n$ and $B$) is the orientation $I_{4}$ (compare \claimref{synthesis.wedge-reflections-compose-to-the-two-pi-rotor-so-geometric-modular-conjugations-give-interacting-spin-statistics}, which writes the Klein twist with $i = I_{4}$). The series does not fix this orientation; (N$_{+}$) assumes it.

\emph{Step 5: one cut and the affine structure, (v).} For $m(u) = e^{i\vartheta(|u|)}$, $m(-u) = m(u)$, so by Step 2's computation the phase of $J_{0}$ is $c\,m(u)^{-1}m(-u) = c$: elementwise geometric with the vacuum's constant phase. The modular group of $\mathfrak{A}((0, \infty))$ for $\omega_{0}\circ\alpha_{m}$ is the vacuum's $\Delta^{is}$ (Step 1), and $\Delta^{is}\pi_{m}(a(f))\Delta^{-is} = \pi_{0}(a(\beta_{s}(mf))) = \pi_{m}(a(m^{-1}\beta_{s}(mf)))$ with $\beta_{s}$ the boost about $0$ (weight $\tfrac12$, rapidity proportional to $s$): $m(u)^{-1}m(e^{s}u) = e^{i(\vartheta(e^{s}|u|) - \vartheta(|u|))}$, which is independent of $u$ for every $s$ exactly when $h = e^{i\vartheta}$ satisfies $h(\lambda x) = \chi(\lambda)h(x)$ for almost every $x > 0$; then $\chi$ is a measurable character of the multiplicative group, $\chi(\lambda) = \lambda^{i\lambda_{0}}$, and $h(x)x^{-i\lambda_{0}}$ is dilation invariant, hence constant: $h(x) = e^{i\mu}x^{i\lambda_{0}}$. So at one cut the elementwise conjugation carries no information on the modular group beyond the vacuum's. The slope-jump vacuum is $\omega_{0}$ transported by the weight-$\tfrac12$ unitary of a regular increasing $\lambda : \R \to \R$ whose slopes at the two ends differ, $\lambda'(u) \to s_{\pm}$ as $u \to \pm\infty$, $s_{-} \neq s_{+}$ (\claimref{synthesis.normality-on-the-closed-generator-gives-the-vacuum-folium-not-the-ground-state}(vi)); its algebra of $O$ is the vacuum's of $\lambda(O)$, so its $J_{a}$ maps the algebra of $O$ to that of $\lambda^{-1}r_{\lambda(a)}\lambda(O)$. This is $r_{a}O$ for every interval $O \subset (a, \infty)$ only if $\lambda(2a - u) = 2\lambda(a) - \lambda(u)$ for all $u > a$; differentiating, $\lambda'(2a - u) = \lambda'(u)$, and letting $u \to \infty$ gives $s_{-} = s_{+}$, a contradiction, at every cut $a$. So at every cut it acts by its own parameter's reflection and by no affine one.

\emph{Step 6: the finite formula, (vi).} On $n_{0} + n_{1}$ modes let $\Omega = a^{*}(\phi_{1})\cdots a^{*}(\phi_{r})\ket{0}$, $E = \sum\ket{\phi_{j}}\bra{\phi_{j}}$, $M$ the CAR algebra of the first $n_{0}$ modes with $\Omega$ cyclic and separating, $X$ the block of $E$ from the first $n_{0}$ modes to the others, $W$ its polar part, $\Gamma_{\Omega} = (-1)^{r}\Gamma$ the parity normalised by $\Gamma_{\Omega}\Omega = \Omega$, and $J_{c}$ the antilinear map $x\Omega \mapsto \varphi(x)\Omega$ for the antilinear multiplicative $*$-map $\varphi$ with $\varphi(a(e_{j})) = \Gamma_{\Omega}\,a^{*}(We_{j})$ (cyclic and separating on the full space forces $n_{0} = n_{1} = r$, so $(-1)^{r} = (-1)^{n_{0}}$ here). The exact check verifies, for $1 + 1$ modes ($\phi = \tfrac35e_{0} + \tfrac{4i}{5}e_{1}$, $W = i$) and $2 + 2$ modes (two Schmidt pairs rotated in each block, $W = i\ket{l_{a}}\bra{r_{a}} + \ket{l_{b}}\bra{r_{b}}$, with $X = W|X|$ checked): the two-point function equals $E$; $\Gamma_{\Omega}\Omega = \Omega$; $J_{c}$ is antiunitary, an involution, and fixes $\Omega$; $J_{c}S$ ($S$ the Tomita operator on the cyclic basis of monomials) is Hermitian with positive leading principal minors, hence positive definite. By uniqueness of the polar decomposition $S = J\Delta^{1/2}$ (Theorem \claimref{singularity.modular-theory-matrix-factor-proved-outright} for the finite case) $J = J_{c}$ and $\Delta^{1/2} = J_{c}S$, and $J_{c}a(e_{j})J_{c} = \varphi(a(e_{j}))$ is checked as an operator identity. The general finite case (factorising the Slater state over Schmidt pairs, each a $1 + 1$ case) is argued, not checked.
